\documentclass[a4paper,10pt]{article}

\usepackage[utf8]{inputenc}
%\usepackage[T1]{fontenc}

\usepackage{textcomp}           % Extra Symbole (Grad Celsius etc.)
\usepackage{amssymb,amsmath}    % Schöne Formeln (AMS = American Mathematical Society)
\usepackage{graphicx}           % Bilder und Seitenränder
\usepackage{subcaption}			% captions for subfigures
\usepackage{booktabs}           % Schönere Tabellen
\usepackage{colortbl}           % Farbige Tabellen

%\usepackage{tcolorbox}			% schöne bunte Boxen
\usepackage{mathtools}			% \mathclap für ordentliche \underbrace-			environments
\usepackage[left=2cm,right=2cm,top=2cm,bottom=2cm]{geometry}			% Pagelayout mit \newgeometry, \restoregeometry
\usepackage{float}
\usepackage{wrapfig}
\usepackage{enumitem}
\usepackage{float}
\usepackage{braket}
\usepackage{caption}
\usepackage{verbatim}
\usepackage[per-mode=reciprocal,output-decimal-marker={.},binary-units=true,separate-uncertainty=true]{siunitx}
\usepackage[breaklinks=true,colorlinks=true,linkcolor=blue,urlcolor=blue,citecolor=blue]{hyperref}
\usepackage{physics}
\usepackage{url}
\usepackage{subcaption}
\usepackage{calrsfs}
\usepackage{tikz}
\usetikzlibrary{decorations, positioning, intersections, calc, shapes,arrows, scopes}
\DeclareMathAlphabet{\pazocal}{OMS}{zplm}{m}{n}

\graphicspath{{./img/}}

\newcommand{\dif}{\mathrm{d}}

\bibliographystyle{instant}

\renewcommand{\k}{\mathbf{k}}
\begin{document}
\begin{titlepage}
 \begin{center}
 \Large{Lab Course in Experimental Physics: Optics}
 \end{center}
 \begin{center}
  \LARGE{\textbf{Focus locking of a tunable lens for optical transportation of ultracold atoms}}
 \end{center}
 \begin{center}
 \large Marco \textsc{Canteri} \\
 marco.canteri@student.uibk.ac.at\\
 \end{center}

 \begin{center}
 \vspace{1cm}
 Innsbruck, \today
 \vspace{1cm}
 \end{center}

 \begin{abstract}
In this project a setup for optical transportation of ultracold atoms has been realized. The focus of a laser was shifted by a tunable optotune lens, controlled by custom electronics and a feedback signal generated from a quadrant photodiode.
\end{abstract}
  \vspace{1cm}

 \begin{center}
 \includegraphics[scale=0.567]{img/uibk}
 \end{center}

\end{titlepage}

\section{Introduction}


\section{Conclusion}

% \begin{thebibliography}{99}
% \bibitem{datasheet_pfeiffer}
% \textsc{Pfeiffer Vakuum}, \textit{Compact FullRange Gauge -PKR 251 }, \url{http://www.idealvac.com/files/brochures/Pfeiffer_PKR_251_Pirani_ColdCathode.pdf}
%
% \bibitem{demtroeder}
% \textsc{W. Demtröder, H.-J. Foth}, \textit{Molekülspektren in kalten Düsenstrahlen}, \url{https://onlinelibrary.wiley.com/doi/pdf/10.1002/phbl.19870430104}
%
% \bibitem{bergmann}
% \textsc{Bergmann-Schäfer}, \textit{Gase, Nanosysteme, Flüssigkeiten}, \textit{Band 5} (de Gruyter, 2005)
%
% \bibitem{cold_cathode}
% \textsc{LDS Vacuum Shopper}, \textit{Cold Cathode Gauges}, \url{http://www.ldsvacuumshopper.com/cocaga.html}
%
% \bibitem{script}
% \textsc{Jennifer Meyer, Roland Wester}, \textit{Versuch im Fortgeschrittenen Praktikum FP3: Supersonische Düsenstrahlen}
%
% \bibitem{lifetime}
% \textsc{Cullen P.J., Milosavljevic V.} (2015) Prog. Theor. Exp. Phys.2015/6, 063J01 . doi:10.1093/ptep/ptv070
%
% \bibitem{illinois}
% \textsc{McCall research group}, \textit{Supersonic Expansion Discharge Source}, \url{http://bjm.scs.illinois.edu/past/source.php}
% \end{thebibliography}
\end{document}
